\section{Discussion and Scope}\label{sec:discussion}

The paper's claim can now be stated compactly. Once a substrate, a packaging, and a staging scale are fixed, observable packaged randomness decomposes into two parts: intrinsic substrate uncertainty and closure deficit. The first term is what remains unpredictable even with exact access to the current microstate. The second is the information about the next packaged state that remains hidden inside the current one because the layer has not stabilized all the distinctions that still matter for prediction. In the controlled Markov benchmark, this deficit vanishes on closed partitions and rises under controlled violations; in budgeted prediction, it is partially bought back by extra memory; and in hashing, it appears as feasibility-limited irreversibility rather than entropy creation. Across these domains, the same explanatory pattern repeats: what a layer calls noise is often unresolved predictive structure.

One advantage of this account is that it separates questions that are often blurred together. Closure deficit is not novelty. A layer may have $\mathsf{CD}_\tau(\Pi)>0$ and still admit a perfectly serviceable ontology; the deficit only says that the current package is not sufficient for one staged step of prediction. Nor is closure deficit the same as genuine irreversibility. Path-space asymmetry and entropy-production audits obey different monotonicity constraints under coarse observation, whereas route mismatch and predictive gaps only register unresolved dependence \citep{tsiokos2026sixbirds,esposito2012stochasticthermodynamicsundercoarsegraining}. Hidden protocol phases, omitted variables, or cheap lifts can therefore make a layer look random without certifying any genuine nonequilibrium drive. The present paper is about the predictive residue of non-closure, not about novelty or arrows of time.

A second advantage is that the account makes randomness explicitly budget-relative without reducing it to a vague appeal to ignorance. Budget matters because records and distinctions are costly to maintain. The curve in \Cref{fig:budget-curve} is therefore not merely a machine-learning performance plot. It is a ledger: increasing memory order lowers the best achievable loss only when the current package has failed to absorb distinctions that still matter for its future. In that sense, randomness is what prediction looks like when closure has been budgeted too cheaply.

The hashing demonstration sharpens the same lesson in an engineered setting. A one-way hash is a deliberately extreme packaging: huge preimage fibers are compressed into a short stable record. But the experiment shows that the appearance of randomness depends as much on the source as on the map. High-entropy inputs produce the familiar small-success $q/2^n$ inversion law and pass simple random-lookingness checks; low-entropy inputs collapse both effects immediately. One-wayness is thus not metaphysical irreversibility. It is packaging observed through a bounded budget, supplied with enough entropy that unpaid distinctions remain expensive to reconstruct.

The auxiliary clustering experiment in \Cref{app:aux-rep} is useful precisely because it does \emph{not} support a simplistic monotonic story. Increasing the number of learned clusters there does not reduce the empirical closure deficit; in the frozen run it increases from $0.0782$ at $k=2$ to $0.1579$ at $k=8$, while held-out packaged prediction worsens at the same time. The lesson is not that the account fails, but that representational budget is not the same as good packaging. More bins do not help when the learned partition is misaligned with the dynamics. Six Birds does not say that refinement is always better; it says that refinement only pays when it closes the right distinctions at the right staging scale.

Several limitations should be kept explicit. The exact treatment here is finite-state and stationary. That is sufficient for the controlled claims we make, but continuous-state or strongly nonstationary settings will require estimators and approximation schemes rather than exact sums. The route-mismatch bound in \Cref{prop:rm-lower-bound} is stated for the stationary-conditioned lift, while many operational pipelines use cheaper lifts for practical reasons. The predictive-gap proxy likewise detects unresolved dependence in packaged history, not literally the micro-closure deficit. These are not objections to the framework; they are reminders that diagnostics should be interpreted as shadows of the exact object rather than substitutes for it.

What the paper does provide is a portable operational answer to a question that is usually posed metaphysically. Randomness need not be treated as a primitive ingredient of explanation at the layer where it is observed. Once the substrate, packaging, staging, and budget are specified, the non-intrinsic part of that randomness is measurable as the information about the next packaged state that remains hidden inside the current one. You usually do not remove such randomness outright; you relocate it, price it, or decide not to pay for it.
